\section{Discussion}
\label{sec:discussion}

Our results show that \textsc{Kaiburr} stays computationally practical despite higher overall security. More specifically,
on both the i7-8700K and the i9-13900K, every \textsc{Kaiburr} parameter set beats FrodoKEM-640 by a wide margin.
Kaiburr-6144, with 1717 classical security, 1557 quantum security, and similar communication costs as FrodoKEM-640, is 
still roughly three times faster on the 8700K (and roughly two-and-a-half times faster on
the 13900K) across KeyGen, Enc, and Dec. In other words, \textsc{Kaiburr}'s \emph{worst} case is still comfortably cheaper than
FrodoKEM's \emph{best} case.

HQC tells an interesting story. Its advantage is concentrated almost entirely in KeyGen: HQC-1, HQC-3, and HQC-5 all beat their corresponding 
\textsc{Kaiburr} parameter set's KeyGen by a wide margin on both machines. That advantage narrows in Enc: HQC still wins clearly against kaiburr-4608 and 
kaiburr-6144, but at the smallest comparison, kaiburr-1792's Enc is at least as fast as HQC-1's Enc on both machines (and noticeably faster for the C/asm variant). 
By Dec, the advantage is reversed: kaiburr-1792's Dec beats HQC-1's Dec by more than a factor of two on both machines, and kaiburr-6144's Dec beats 
HQC-5's Dec as well. Only the middle parameter set, kaiburr-4608 and HQC-3, have similar Dec performance.

The ACC results show how the cost breakdown changes as $k$ increases. \texttt{poly\_gen\_matrix} accounts for 73--84\% of total cycles
(Table~\ref{tab:acc-profile}), and its cost grows quadratically with $k$ since a $k \times k$ matrix requires $k^2$ sampled polynomials.
Hardware SHAKE-128 acceleration does not remove this bottleneck either. The core waits on the accelerator for only 5.5--6.2\% of total cycles, with little change 
across parameter sets (Table~\ref{tab:acc-keccak-all}). Instead, rejection sampling accounts for an increasing share of the
total cost, reaching 69.9\% for kaiburr-6144. Our analysis suggests that further hardware optimization should focus on rejection sampling 
rather than SHAKE throughput. This is consistent with prior work by Abdulrahman et al., who also identify rejection sampling as a 
target for further hardware optimization~\cite{SP:AOPPSS25}. In our setting, its cost grows with $k$ and becomes
the dominant part of matrix generation.

As discussed in Section~\ref{sec:introduction}, our work explicitly explores increasing the security margin as a hedge against improvements in the best known 
attacks today. \textsc{Kaiburr} shows that you can buy enormous security margin without abandoning the efficiency advantages of module structure, while requiring 
much less computational cost than competitive unstructured schemes. This provides conservative security in the sense of margin, and not the assumption underlying 
the scheme. By this, we mean that \textsc{Kaiburr} still is and will remain vulnerable to \emph{structural} attacks that exploit the module structure, 
a class of attacks FrodoKEM is immune against. The purpose of designing and evaluating these high security variants of ML-KEM (Table~\ref{tab:kaiburr-params}) 
is not to argue that such security margins are currently necessary. Considering the relative stability of lattice cryptanalysis and security estimates in 
recent years, we do not expect improvements of the magnitude that would require such parameter sets in the near future. Our work is focused on 
characterizing how much margin ML-KEM can accommodate and the resulting engineering tradeoffs. Cryptographically, the natural response to such improvements 
in attack methods and changing security estimates has been to scale parameters. We have shown that applying this naive approach to ML-KEM eventually encounters 
a correctness limit (Table~\ref{tab:cbd-baseline}) and scaling is unavailable without changing the noise distribution. \textsc{Kaiburr} also shows that 
this limit can be pushed substantially further while retaining the \emph{essence} of ML-KEM that gives it its efficiency advantages. As a result, 
instead of treating the margin as a deployment requirement, we study the tradeoff between security margin and computational cost to understand the consequences 
of scaling, if it ever becomes necessary.